% Part: sets-functions-relations
% Chapter: functions
% Section: isomorphisms

\documentclass[../../../include/open-logic-section]{subfiles}


\begin{document}

\olfileid{sfr}{fun}{iso}
\olsection{সমরূপতা}

\begin{explain}
একটি \emph{সমরূপতা} হল এমন এক-এক ও সর্বব্যাপী চিত্রণ, যা তার দ্বারা
সম্পর্কিত সেটগুলির গঠন রক্ষা করে। এখানে গঠন বলতে সেটগুলির
!!{element}s-এর মধ্যে বিদ্যমান সম্পর্কগুলি বোঝায়।
দুটি সেট $X=\{1,2,3\}$ ও $Y=\{4,5,6\}$ বিবেচনা করো।
উভয় সেটেই উত্তরসূরি, ছোট হওয়া ও বড় হওয়ার সম্পর্কগুলি গঠন দেয়।
সেটদুটির মধ্যে একটি সমরূপতা হল সেই গঠনগুলি রক্ষাকারী !!a{bijection}।
সুতরাং একটি !!a{bijective} অপেক্ষক $f \colon X \to Y$
সমরূপতা হয়, যদি নিচের শর্তগুলি সিদ্ধ হয়: $i<j$ যদি এবং কেবল যদি $f(i)<f(j)$,
$i>j$ যদি এবং কেবল যদি $f(i)>f(j)$, এবং
$j$ হল $i$-এর উত্তরসূরি যদি এবং কেবল যদি $f(j)$ হল
$f(i)$-এর উত্তরসূরি।
\end{explain}

\begin{defn}[সমরূপতা]
ধরা যাক $U$ হল $\langle X, R\rangle$ জোড়া এবং
$V$ হল $\langle Y, S\rangle$ জোড়া, যেখানে $X$ ও $Y$ সেট,
আর $R$ ও $S$ যথাক্রমে $X$ ও $Y$-র উপর সম্পর্ক।
$X$ থেকে $Y$-তে একটি !!{bijection} $f$ হল $U$ থেকে $V$-তে
\emph{সমরূপতা}, যদি এবং কেবল যদি সেটি সম্পর্কীয় গঠন রক্ষা করে।
অর্থাৎ, $X$-এর যেকোনো $x_{1}$ ও $x_{2}$-র জন্য
$\tuple{x_1,x_2} \in R$ যদি এবং কেবল যদি
$\tuple{f(x_1),f(x_2)} \in S$।
\end{defn}

\begin{ex}
দুটি সেট $X=\{1,2,3\}$ ও $Y=\{4,5,6\}$ এবং ছোট হওয়া ও
বড় হওয়ার সম্পর্কগুলি বিবেচনা করো। $f(x) = 7-x$ দ্বারা সংজ্ঞায়িত
অপেক্ষক $f\colon X \to Y$ হল $\tuple{X,<}$ ও $\tuple{Y,>}$-র
মধ্যে একটি সমরূপতা।
\end{ex}
% BN-SRC-910 TARGET-CORRECTION-BEGIN
% Bengali conditional grammar repaired; all relation-preservation clauses unchanged.
% BN-SRC-910 TARGET-CORRECTION-END

\end{document}
